% TftModule
\section{【中級】TftModule — SPI共有 + 複合デバイスの分離}

\begin{patternbox}{TftModule で学べるパターン}
\begin{itemize}
    \item SPI共有ドライバを使った画面描画
    \item ModuleTimerによるFPS制御（更新頻度の制限）
    \item テキスト表示をData構造体経由で受け渡す方法
\end{itemize}
\end{patternbox}

\subsection{Data構造体のテキストバッファ}

\begin{lstlisting}[style=cppstyle, caption={TftModule.h — Data}]
struct TftData {
    bool touchPressed = false;
    int16_t touchX = 0, touchY = 0;
    char line1[48] = "";  // ロジックフェーズが書き込む表示テキスト
    char line2[48] = "";
    char line3[48] = "";
    char line4[48] = "";
    char line5[48] = "";
};
\end{lstlisting}

\begin{pointbox}[出力モジュールは「何を表示するか」を決めない]
\texttt{line1}〜\texttt{line5}はロジックフェーズが\texttt{snprintf()}で書き込む文字列。
TftModuleの\texttt{update()}はそれを画面に描画するだけで、
表示内容の判断は一切行わない。出力モジュールの責務を守っている。
\end{pointbox}

\subsection{ロジックフェーズでの表示内容設定}

\begin{lstlisting}[style=cppstyle, caption={main.cpp — applyPattern()でのTFT表示設定}]
void applyPattern(SystemData& data) {
    // MPU-6500データをTFT表示用にフォーマット
    if (data.mpu.isValid) {
        snprintf(data.tft.line1, sizeof(data.tft.line1),
                 "Ax:%.2f Ay:%.2f Az:%.2f",
                 data.mpu.accelX, data.mpu.accelY, data.mpu.accelZ);
    }
    // カメラ情報
    if (data.camera.isValid && data.camera.frameReady) {
        snprintf(data.tft.line4, sizeof(data.tft.line4),
                 "CAM: %dx%d (%d B)",
                 data.camera.width, data.camera.height,
                 (int)data.camera.frameSize);
    }
}
\end{lstlisting}

\begin{pointbox}[入力→ロジック→出力のデータフロー]
入力フェーズでMpu6500Moduleが\texttt{data.mpu}に値を書き込み、
ロジックフェーズでそれを文字列に変換して\texttt{data.tft.line1}に書き込み、
出力フェーズでTftModuleが画面に描画する。
3フェーズモデルの典型的なデータフローの実例。
\end{pointbox}
